\section{Robustness}
\label{sec:robustness}

The mortality null is the paper's identified result, so we stress it from
every direction we can. It survives variation in the closure sample, the
weighting, the estimator, the pandemic window, and the exposure
threshold. We report the checks in turn; the headline---no detectable
increase in suicide or self-harm mortality---does not move.

\paragraph{§6.1 — Closure sample.} The primary sample is the $\valNclosuresPnash$
PNASH-anchored psychiatric closures. Table~\ref{tab:mortality-null}
(bottom panel) re-estimates the population-weighted suicide effect on two
alternative closure sets: the $\valNspecClosures$ specialized closures surviving the
full statistical filter ($\valSuicSpecAtt$, 95\% CI $[\valSuicSpecLo, \valSuicSpecHi]$) and the
broader set of $\valNclosuresPsymax$ psychiatric closures formed by relaxing the
demand-decline filter ($\valSuicPsymaxAtt$, $[\valSuicPsymaxLo, \valSuicPsymaxHi]$). Both are null and
both rule out increases above roughly one-quarter of baseline. The null
is not an artifact of how the closure sample is drawn.

\paragraph{§6.2 — Weighting.} Population- and municipality-weighted
estimates appear side by side in Table~\ref{tab:mortality-null}. The
population-weighted suicide effect is $\valSuicWAtt$; the municipality-weighted
estimate is $\valSuicOAtt$. They differ in sign in their (statistically
indistinguishable from zero) point estimates because the effect, to the
extent there is one, differs between small and large municipalities, but
both confidence intervals exclude large increases. We lead with the
population-weighted estimate because it is the policy-relevant quantity
(harm per person), and report the unweighted estimate so the null is
visibly not a weighting artifact.

\paragraph{§6.3 — Estimator.} Under the heterogeneity-robust
Sun-Abraham estimator the suicide effect is $\valSuicWAtt$ (SE $\valSuicWSe$). The
naive two-way fixed-effects estimator returns a larger $\valSuicTwfeAtt$ (SE
$\valSuicTwfeSe$)---but two-way fixed effects is biased under treatment-effect
heterogeneity across cohorts \citep{deChaisemartin2020, sunAbraham2021},
exactly the contamination the modern estimators correct, and the
heterogeneity-robust estimators shrink it to zero. Even the contaminated
estimator yields only a modest effect that the credible estimators
null out; there is no large mortality response hiding under any
specification. The first stage and the mortality null are likewise
reproduced under \citet{cs2021} and \citet{borusyak2024}.

\paragraph{§6.4 — Pandemic window.} The COVID-$19$ years $2020$--$2021$
disrupted admissions, mobility, and cause-of-death coding. Dropping them
leaves the suicide effect at $\valSuicDropPandAtt$ (SE $\valSuicDropPandSe$), within a tenth of the
main estimate. The pandemic is not driving the null.

\paragraph{§6.5 — Exposure threshold.} The flow exposure rule sets the
share cutoff at $\theta = 0.05$. The first-stage travel-burden result is
negative and significant across the continuous sweep
$\theta \in [0.01, 0.12]$, confirming that the exposed set is not an
artifact of the default cutoff; the mortality null holds across the same
range, since a null estimated on a larger or smaller exposed set remains
null. Full sweep diagnostics are in Appendix~\ref{app:theta}.

\paragraph{§6.6 — Spillover and state capitals.} If a closure redirects
patients to a hospital also serving an untreated municipality, the
control could be contaminated. Excluding the $27$ state capitals---the
principal referral hubs and the main sites of such congestion---leaves
the suicide effect at $\valSuicNoCapAtt$ (SE $\valSuicNoCapSe$), if anything a tighter null,
and changes the first-stage travel ATT by less than $1\%$. The result is
not driven by the largest referral destinations.

\paragraph{§6.7 — Inference with few effective clusters.} The panel holds
thousands of municipalities, but identification rests on
$\valNclosuresPnash$ closures, so municipality-clustered standard errors
could overstate precision if shocks are correlated within a closure's
catchment. Two checks address this. Clustering at the Federation Unit
level ($\valNUFclusters$ clusters) widens the suicide standard error from
$\valSuicWSe$ to $\valSuicSeUF$ and the confidence interval to
$[\valSuicUFLo, \valSuicUFHi]$---still centered near zero and still
excluding the large increases critics predicted. Randomization inference
is more conservative still: permuting the cohort treatment assignment
across municipalities yields a placebo distribution against which the
observed effect sits at a two-sided $p = \valSuicRIp$, well inside the
placebo range $[\valSuicRILo, \valSuicRIHi]$. Neither check overturns the
null, and the design-based placebo spread still rules out increases above
roughly a quarter of baseline---evidence that the result reflects an
absence of large effects, not understated uncertainty.

\paragraph{Summary.} Across closure samples, weightings, estimators, the
pandemic window, the exposure threshold, spillover controls, and the
inference procedure, the mortality null is stable. The one place the data move---the naive TWFE
estimator---is the place heterogeneity-robust methods are designed to
correct, and they do. We read the pattern as evidence that the null
reflects the absence of a large mortality response, not a fragile
specification. The parallel first-stage and ICSAP robustness battery
(eighteen specifications) is reported in Appendix~\ref{app:robustness}.
